% !TEX root = ../main.tex

\chapter{Design Specification}

\section{Customer Research Method}

\subsection{Research Objective and Instrument Design}

The primary stakeholder is a researcher who needs to triage new papers, inspect a paper's contribution, ask evidence-seeking questions, revisit related work, and refine future discovery.
The supplied materials also imply secondary stakeholders: project teams that may share papers, system operators who manage model and storage costs, and paper authors whose work must not be misrepresented.
The customer study investigated how these users currently discover, select, store, recover, connect, and verify papers.
It was deliberately conducted before presenting MNEME so that participants would describe existing behavior rather than confirm a proposed product.

The interview guide contained eleven open-ended questions.
The sequence moved from research context and reading frequency, through current tools and recent discovery episodes, to overload, organization, recall, conceptual connections, and trust in AI-produced summaries.
Behavior-grounded prompts such as ``tell me about the last time'' were used to obtain concrete episodes.
Follow-up prompts asked what happened next, what made the episode difficult, and how it affected the participant's work.
This funnel design reduced the risk that a proposed feature would be embedded in the question itself.
The full instrument is retained in Appendix~\ref{app:interview}.

\subsection{Participants, Procedure, and Evidence Boundary}

The project report records ten semi-structured interviews, coded P1--P10.
Participants ranged from first-year master's students to doctoral researchers and a lecturer, and represented machine learning/NLP, computer vision, and program verification.
They were recruited through laboratories, paper-reading groups, and campus study areas.
Sessions were conducted individually in person or by video call and lasted approximately fifteen minutes on average.
All four team members contributed to instrument design, recruitment, or interviewing.
Notes were transcribed immediately after each session into single-observation statements for affinity mapping.

This sample is appropriate for formative problem framing but not for population inference.
It is small, non-random, and concentrated in computing-related disciplines.
The local evidence package contains the consolidated report and participant-coded affinity map, but not full transcripts, audio, consent forms, or a complete coding audit trail.
Accordingly, counts in this chapter indicate how the team organized the ten interviews.
They are not prevalence estimates for academic researchers as a whole.

\section{Interview Findings and Affinity Synthesis}

\subsection{Affinity-Mapping Procedure}

Each reported observation was assigned to one digital sticky note and tagged with its participant code.
Notes were clustered bottom-up into subthemes and then into four top-level groups: Pain Points, Current Solutions, Trust in AI, and Wishes.
The report lists 36 pain-point notes, 74 current-solution notes, 21 trust-related notes, and 2 wish notes.
These totals describe coded observation units rather than unique people, so they must not be interpreted as percentages.
Their value is structural: the map preserves the relationship between individual observations, repeated sentiments, and higher-level design themes.

\begin{figure}[htbp]
  \centering
  \includegraphics[page=3,width=0.93\textwidth]{figures/source_materials/design_specification.pdf}
  \caption{Reported affinity map and coding explanation from the customer-research package. The embedded page is retained as a source artifact; the discussion in the text preserves its small-sample limitation.}
  \label{fig:affinity-map}
\end{figure}

\subsection{Pain Points and Existing Workarounds}

The largest recurring problem categories were information overload, difficulty retrieving a previously encountered paper, storage decay, and the absence of a maintained personal map of relationships.
The report records episodes in which a participant remembered an idea but not a searchable title or author, spent tens of minutes recovering one paper, or asked a colleague to identify it.
Other observations concerned folders or topic lists that became ``black holes'' because inconsistent naming and blank metadata weakened later search.
These episodes support a downstream framing: search is necessary, but it does not guarantee durable comprehension or recall.

Participants already used a heterogeneous tool chain.
Discovery channels included Google Scholar, arXiv, advisor or colleague referrals, mailing lists, and general-purpose AI tools.
Management practices included Zotero, Mendeley, local folders, browser bookmarks, spreadsheets, and note systems.
Reading decisions used title, abstract, author, venue, citation context, recency, and recommendation by trusted people.
The design implication is not to replace every existing tool.
MNEME should preserve original links and bibliographic identity, minimize migration cost, and concentrate on the transitions that existing workflows leave to manual effort.

\subsection{Trust and Verification}

Trust-related observations converged on inspectability.
Participants wanted to verify important statements against the paper and were wary of confident language without visible support.
Some already used AI assistance, but described manual checking as necessary.
This finding shaped two requirements.
First, summaries and answers must expose source actions near the relevant claim rather than in a remote bibliography.
Second, the interface must distinguish missing evidence, partial evidence, and verified support instead of collapsing them into a binary badge.

\section{Customer Profile and Value Proposition}

The synthesized customer profile is an early-stage or active academic researcher who follows a changing literature, already has entrenched search and storage practices, and needs a lower-friction way to decide, understand, recover, and connect papers.
The customer's ``job'' is not merely to collect PDFs.
It is to maintain a defensible working understanding of a field while moving between short triage sessions and deeper reading.
The profile therefore values relevance, traceability, low setup cost, compatibility with existing sources, and the ability to correct an automated interpretation of interest.

\begin{figure}[htbp]
  \centering
  \includegraphics[page=4,width=0.93\textwidth]{figures/source_materials/design_specification.pdf}
  \caption{Customer-profile synthesis in the supplied design specification. It connects reported jobs, pains, current tools, and trust expectations.}
  \label{fig:customer-profile}
\end{figure}

The value proposition is summarized as \emph{read, remember, and plan with evidence}.
MNEME reduces an unbounded paper stream to a limited briefing, provides concise explanations that retain a source path, records explicit and bounded behavioral signals, and exposes a small related-paper neighborhood.
The proposition differs from passive search because it continues after discovery, and differs from an autonomous agent because every consequential transition is intended to remain inspectable and correctable.

\section{Workflow Needs and Design Principles}

Six design principles convert the customer synthesis into constraints.
\emph{Continuity} requires one identity across digest, paper detail, question answering, local cache, and graph.
\emph{Evidence proximity} keeps verification actions adjacent to generated claims.
\emph{Progressive disclosure} shows a brief explanation first and permits deeper source inspection without overwhelming the mobile screen.
\emph{Correctability} lets users edit interests and reject misleading adaptations.
\emph{Graceful degradation} preserves a useful abstract-level or cached state when parsing, connectivity, or providers fail.
\emph{Bounded attention} limits briefing size, notification frequency, graph size, and background work.

The intended workflow begins with a digest.
Each recommended paper includes metadata, a concise summary, and a recommendation reason.
Opening the paper reveals a structured summary and evidence controls.
A question opens a scoped retrieval-and-generation interaction.
The user may then inspect citation relations or modify interest settings.
Events generated by these actions can inform later recommendations, but only according to the event semantics defined in Chapter 3.

\begin{figure}[htbp]
  \centering
  \includegraphics[page=4,width=0.93\textwidth]{figures/source_materials/storymap_architecture.pdf}
  \caption{Five-stage story map used to distinguish skeletal, MVP, and stretch scope. The thesis maps these product stages to TC1--TC4 rather than treating every feature as a research contribution.}
  \label{fig:story-map}
\end{figure}

\section{Competitor Categories and Unmet Design Space}

The local competitor analysis groups existing options into discovery engines, reference managers, AI paper assistants, and general chat systems.
Discovery engines provide broad and familiar search but do not maintain an evidence-linked personal reading loop.
Reference managers provide durable bibliographic organization but depend on manual tagging, note taking, and retrieval habits.
AI paper assistants reduce the cost of summarization and question answering but vary in source transparency and do not necessarily preserve exact-revision lineage or a user-correctable interest state.
General chat systems are flexible but place even more responsibility on the user to define scope and verify answers.

MNEME's intended position is therefore combinational rather than feature-exclusive.
Its differentiation depends on joining proactive delivery, source-linked assistance, behavior-aware but inspectable personalization, and mobile graph exploration.
This position must be evaluated as a loop: if any one of provenance, evidence support, interest correction, or mobile action fails, the combined claim narrows.
Competitor facts that depend on current commercial product behavior should be dated and rechecked before final submission; the thesis relies primarily on category-level comparison and the directly observed customer workflow.

\section{Functional Requirements}

Table~\ref{tab:functional-requirements} condenses the feature plan into research-facing requirements.
The requirements are grouped by the invariant they protect rather than by implementation library.

\begin{longtable}{P{1.0cm}P{4.7cm}P{5.0cm}P{2.8cm}}
\caption{Functional requirements mapped to the thesis challenges.}
\label{tab:functional-requirements}\\
\toprule
ID & Requirement & Acceptance evidence & Traceability \\
\midrule
\endfirsthead
\toprule
ID & Requirement & Acceptance evidence & Traceability \\
\midrule
\endhead
R1.1 & Assign a stable logical identity and an explicit revision identity to every ingested paper. & Version fixtures show that two revisions remain distinguishable while belonging to one logical paper. & TC1--E1 \\
R1.2 & Propagate revision and transformation metadata to sections, chunks, summaries, embeddings, and answer citations. & A provenance query reconstructs every dependency for a displayed claim. & TC1--E1 \\
R1.3 & Label parsing quality and retain a defined degraded mode. & Multi-layout fixtures produce full, partial, or abstract-only states without silent success. & TC1--E1 \\
R2.1 & Scope retrieval by paper, selected library, or declared collection. & Retrieval traces contain only documents allowed by the selected scope. & TC2--E2 \\
R2.2 & Represent generated answers as claims linked to source spans. & Every externally checkable claim contains resolvable evidence or an explicit unsupported state. & TC2--E2 \\
R2.3 & Distinguish identifier, availability, consistency, and coverage checks. & Failure fixtures trigger the intended verification label. & TC2--E2 \\
R2.4 & Abstain when available evidence does not support an answer. & Unanswerable questions are evaluated separately from answerable questions. & TC2--E2 \\
R3.1 & Record explicit and implicit events with documented semantics and deduplication identifiers. & Replayed event batches are idempotent and event provenance is inspectable. & TC3--E3 \\
R3.2 & Maintain one bounded, time-decayed normalized behavior-v1 vector; keep separable multi-timescale state as planned behavior-v2. & Synthetic event replays yield predictable bounded and decayed behavior-v1 updates. & TC3--E3 \\
R3.3 & Explain recommendation reasons using inspectable evidence from the interest state. & Reasons map to declared topics, authors, papers, or events. & TC3--E3 \\
R3.4 & Permit users to edit, suppress, or delete inferred interests and behavioral history. & Correction tests confirm that the visible and ranking states change consistently. & TC3--E3 \\
R4.1 & Preserve essential reading state during intermittent connectivity. & Offline and reconnection scenarios retain consistent local state. & TC4--E4 \\
R4.2 & Bound background work, push frequency, and model cost. & Configuration and logs demonstrate rate, frequency, and budget limits. & TC4--E4 \\
R4.3 & Expose degraded states and verification actions in the mobile UI. & Usability tasks measure whether users can identify evidence and recover from failures. & TC4--E4 \\
\bottomrule
\end{longtable}

\section{Scope Boundaries and Acceptance Levels}

The story map separates three maturity levels.
The \emph{skeletal product} demonstrates an end-to-end path: fetch papers, retain basic metadata and an abstract or TLDR, collect explicit interests, assemble a digest, and support a basic single-paper question flow.
The \emph{MVP} adds section-aware preparation, embeddings, behavioral events, explainable ranking, bounded notifications, source-linked summaries, citation checks, and graph exploration.
\emph{Stretch} items include voice input, camera OCR, autonomous multi-hop research, team collaboration, Zotero/BibTeX export, and on-device model research.

This separation prevents an incomplete convenience feature from blocking evaluation of the central research loop.
It also prevents stretch items from being presented as contributions before core evidence exists.
For each feature, completion requires a visible user outcome, a defined failure state, and an executable acceptance check.
Appendix~\ref{app:acceptance} preserves the product-level Given--When--Then criteria, while Chapters 3 and 4 define the lower-level mechanism and evaluation evidence.

\section{Non-Functional and Ethical Requirements}

\subsection{Provenance and Reproducibility}

Every claim-bearing artifact must record enough configuration to be reproduced: source document identity and revision, parser version, chunking configuration, embedding model, retrieval parameters, generator and prompt version, and verification policy.
The requirement does not imply that all third-party models are deterministic.
It requires the system to preserve the conditions under which an output was produced.

\subsection{Privacy and Data Minimization}

Behavioral events are potentially sensitive because they reveal research topics, reading habits, and professional activity.
The event schema should therefore collect only fields required by the declared update model.
Users must be able to inspect and remove events and inferred interests.
Retention, export, deletion, and synchronization policies must be explicit.
The system must not interpret a missing click as a stable negative preference without a stated exposure condition.

\subsection{Uncertainty and User Control}

Human--AI interaction guidelines emphasize communicating capability, supporting correction, and responding to feedback~\cite{amershi2019guidelines}.
For MNEME, uncertainty is not represented by a single confidence number.
The interface should instead expose operational states: evidence retrieved, source span available, consistency check passed or failed, answer incomplete, or no answer supported.
Users should be able to open the source, change retrieval scope, reformulate a question, and correct the interest state.

\subsection{Latency, Reliability, and Cost}

Latency targets must be measured per operation rather than summarized by one end-to-end average.
Digest loading, paper detail retrieval, local navigation, first-token latency, complete-answer latency, graph expansion, and synchronization have different user consequences.
Reliability similarly requires error classes: network failure, provider failure, parse failure, retrieval-empty, verification failure, and local synchronization conflict.
Cost boundaries should include a daily model budget and cached results keyed by content and configuration hashes.

\section{Interaction Design}

\subsection{Digest and Paper Detail}

The interactive prototype presents the digest as the home screen.
A paper card exposes its title, authors, a concise description, a match indicator, a recommendation reason, and an open action.
This organization is intended to support rapid triage.
The paper-detail view then separates summary content from source-verification actions so that the user can move from overview to evidence.
These statements describe the prototype and do not establish production behavior.

\subsection{Cited Question Answering}

The Q\&A view is intended to display an answer together with source chips or evidence spans.
The interaction must not label an answer ``verified'' merely because a paper identifier is valid.
The user-visible status should correspond to the verification layers defined in Chapter 3.
Multi-turn interaction must also preserve the selected scope and make scope changes visible.
Recent multi-turn RAG evaluation shows that follow-up, unanswerable, and non-standalone questions create distinct challenges~\cite{katsis2025mtrag}; therefore, follow-up convenience cannot replace explicit evaluation.

\subsection{Knowledge Graph and Interest Controls}

The graph view presents related or citing papers as nodes.
The interest view permits topic and preference adjustments.
These screens support two different forms of continuity: the graph externalizes relations among papers, while the interest view externalizes part of the system's user model.
Both require clear actions.
A visual node without a visible selected state or primary action can appear decorative rather than interactive.
An interest label without provenance or correction controls can appear authoritative even when inferred from weak evidence.

\section{Formative Prototype Evidence}

The final usability presentation reports that five participants completed five prototype tasks.
It reports success counts of 5/5 for opening a recommended paper, 5/5 for identifying evidence, 5/5 for cited Q\&A, 2/5 for opening a related paper from the graph, and 4/5 for changing preferences.
The presentation interprets graph discoverability as the primary drop-off and proposes an instruction, a selected-node state, and a persistent ``Open paper'' action.

\begin{figure}[htbp]
  \centering
  \includegraphics[page=6,width=0.93\textwidth]{figures/source_materials/usability_tests.pdf}
  \caption{Aggregate task outcomes reported by the formative-usability presentation. These counts motivate redesign but are not treated as final system effectiveness evidence.}
  \label{fig:formative-task-results}
\end{figure}

These counts are classified as \observed{}/\limited{}, not \measured{} final evidence.
The supplied folder does not include the participant-by-task table, timing and tap records, coded observations, consent documentation, or locally accessible recordings.
The companion draft presentation explicitly contains blank fields for those artifacts.
Consequently, the counts can motivate a design iteration, but they cannot establish a general usability effect or validate the redesign.
\section{Mobile Acceptance and Evidence Boundaries}

The mobile requirements become testable when each one names a reader-visible result under a defined condition.
Table~\ref{tab:mobile-acceptance} extends the existing R4 requirements with client-side acceptance conditions.
It specifies what the Android client must expose; Chapter~4 reports the corresponding E4 measurements.

\begin{table}[htbp]
\centering
\caption{Client-side acceptance conditions for bounded mobile integration.}
\label{tab:mobile-acceptance}
\begin{tabular}{P{2.4cm}P{6.2cm}P{5.3cm}}
\toprule
Requirement link & Observable acceptance condition & E4 evidence \\
\midrule
R4.1: reading continuity &
A stored backend briefing remains available when refresh fails, carries a cached-content label, and is never substituted for live content without disclosure.
When no valid stored briefing exists, the client presents a retryable error. &
Controlled live, cached, offline, server-failure, and invalid-cache trials. \\
\addlinespace
R4.1 and R4.3: asynchronous work &
Paper preparation exposes loading, processing, ready, and error states.
Polling terminates at a completed or failed job and does not leave the interface in an unexplained waiting state. &
Controlled three-stage polling and startup/resume measurements. \\
\addlinespace
R4.3: graph interaction &
The graph declares ready, fallback, or empty state.
A selected node remains visible in the native interface and provides an Open Paper action that returns to the same graph selection. &
Graph-state trials at the declared client bound and graph/navigation software tests. \\
\addlinespace
R3.1 and R4.1: event return &
Visible paper interactions enter a durable queue with stable identifiers.
Retryable failures preserve pending work, and replayed identifiers are acknowledged as duplicates. &
Controlled queue/retry/replay trials and a separate live-backend trace. \\
\bottomrule
\end{tabular}
\end{table}

These conditions support implementation and systems acceptance.
They do not establish whether readers understand every status label or whether the graph redesign improves task completion.
No participant retest accompanied the Android measurements, so those user-facing effects remain unestablished.

The implemented event payload also narrows the privacy boundary.
Paper impressions, paper opens, and submitted paper-scoped questions are represented by an event type, relevant identifiers, and a timestamp.
The submitted question text is not copied into the synchronization event.
This data-minimization decision reduces the behavioral content held in the client queue; it is not a complete privacy or security audit.
